Our research intersects with multiple strands of literature focusing on LLM alignment, conversational biases, and the mechanics of reasoning under pressure. We categorize the related work into three primary themes: dialogic deference, sycophancy in causal judgments, and the general biases of LLM-as-a-judge frameworks.

\para{Dialogic Deference and Conversational Bias}
A growing body of work explores how conversational framing influences LLM behavior. \citet{rabbani2026dialdefer} introduce the concept of ``Dialogic Deference,'' demonstrating that models shift their judgments significantly depending on whether a claim is framed as an objective fact or attributed to a conversational speaker. Their findings emphasize that LLMs are sensitive to social dynamics, frequently shifting toward agreement (deference) or disagreement (skepticism) based on the domain. Building on their exploration of conversational framing, we focus specifically on the intensity of adversarial or judgmental tones and quantify their impact on the length and quality of the intermediate reasoning steps, rather than just the final verdict shift.

\para{Sycophancy and Skepticism Traps}
The tendency of LLMs to alter their answers to please users has been well-documented. \citet{chang2026diagnosing} identify the ``Skepticism Trap'' and ``Sycophancy Trap'' in causal reasoning tasks, showing that safety-tuned models may reject valid causal links to avoid over-claiming, and flip correct rejections to endorsements when placed under social pressure. Similarly, \citet{kim2025challenging} investigate sycophancy under user rebuttal, revealing that LLMs frequently endorse a user's counterargument when challenged sequentially. Unlike these studies, which primarily view user pressure as a vulnerability that induces sycophantic failure, we examine the hypothesis that such pressure can also serve as a catalyst for deeper, more rigorous analytical thought---effectively establishing a constructive ``sweet spot'' for skepticism.

\para{Biases in LLM-as-a-Judge}
As LLMs are increasingly utilized to evaluate text and behavior, their inherent biases have come under scrutiny. \citet{ye2024justice} categorize twelve types of biases in LLM-as-a-Judge systems, including authority bias and verbosity bias. Their methodology involves injecting perturbations to measure consistency. We complement this line of research by isolating prompt tone as a primary variable, demonstrating that the verbosity induced by judgmental prompts is not necessarily a superficial bias, but often reflects a substantive increase in task-relevant reasoning that leads to more accurate outcomes in highly nuanced domains.